In the static orthonormal coframe $\theta^{0} = \sqrt f\,dt$, $\theta^{1} = dr/\sqrt f$, $\theta^{2} = r\,d\vartheta$, $\theta^{3} = r\sin\vartheta\,d\varphi$, $f = 1 - 2M/r$, the connection forms are $\omega^{0}{}_{1} = (f'/2\sqrt f)\,\theta^{0}$, $\omega^{2}{}_{1} = (\sqrt f/r)\,\theta^{2}$, $\omega^{3}{}_{1} = (\sqrt f/r)\,\theta^{3}$, $\omega^{3}{}_{2} = (\cot\vartheta/r)\,\theta^{3}$, and the curvature operator is diagonal on the coordinate planes with eigenvalues $2M/r^{3}$ on the radial plane $\hat0\hat1$ and the angular plane $\hat2\hat3$ and $-M/r^{3}$ on the four mixed planes. The extreme eigenvalue distinguishes the pair $\{\sigma_{r}, I\sigma_{r}\}$ invariantly, the radial plane and its Hodge dual; the radial member is selected within the pair by the sign of the square, $\sigma_{r}^{2} = +1$ against $(I\sigma_{r})^{2} = -1$. Which of the two null directions of the radial plane is the vacuum's is the flag bit, which the series takes as initial data. In Painlev\'e-Gullstrand form $ds^{2} = dT^{2} - (dr + \beta\,dT)^{2} - r^{2}d\Omega^{2}$ with $\beta^{2} = 2M/r$, regular through the horizon. \emph{Provenance:} the connection and curvature forms \cite[App.~A.2]{ButlerT5}; the pair, the sign of the square and the matching \cite[Thm.~6.1, Prop.~7.21(3)]{ButlerT5}; the Painlev\'e-Gullstrand form \cite[App.~A.1]{ButlerT5}; the flag bit \cite[Def.~2.10, Prop.~5.1]{ButlerT14}.
